% Part: first-order-logic
% Chapter: introduction
% Section: syntax

\documentclass[../../../include/open-logic-section]{subfiles}

\begin{document}

\olfileid{fol}{int}{syn}

\olsection{সংকেতবিন্যাস}

প্রথমে নির্ভুলভাবে স্থির করতে হবে, কোন কোন সংকেতক্রম প্রথম-ক্রমের
যুক্তিবিদ্যার !!{sentence}s হিসেবে গণ্য হবে। পরে আমরা তা করব; আপাতত
উদাহরণ ধরে এগোব। প্রথম-ক্রমের যুক্তিবিদ্যার মৌলিক গঠন-উপাদান—তার
শব্দভাণ্ডার—দুই ভাগে বিভক্ত। প্রথম ভাগে আছে সেই সংকেতগুলি, যেগুলির
সাহায্যে আমরা নির্দিষ্ট কিছু বলি বা নির্দিষ্ট কোনো কিছুকে নির্দেশ করি।
আমরা !!{constant}s দিয়ে বস্তু নির্দেশ করি এবং নির্দেশিত বস্তুগুলি
সম্বন্ধে !!{predicate}s দিয়ে কিছু বলি। যেমন, কোনো একক বস্তুকে নির্দেশ
করতে !!a{constant} হিসেবে $\Obj a$ ব্যবহার করতে পারি, তারপর
!!{sentence}~$\Atom{\Obj P}{\Obj a}$ দিয়ে বস্তুটি সম্বন্ধে কিছু বলতে
পারি। “$\Obj a$” ও “$\Obj P$”-এর কোনো অর্থ মনে ধরে নিলে
$\Atom{\Obj P}{\Obj a}$-কে ইংরেজির একটি বাক্য হিসেবে পড়তে পারো (এবং
বিধিবদ্ধ যুক্তিবিদ্যা প্রথম শেখার সময় সম্ভবত তা করেছও)। প্রথম-ক্রমের
যুক্তিবিদ্যার এমন সরল !!{sentence}s পাওয়ার পর শব্দভাণ্ডারের দ্বিতীয়
ভাগ—যৌক্তিক সংকেতগুলি (সংযোজক ও পরিমাণসূচক)—ব্যবহার করে আরও জটিল বাক্য
গঠন করা যায়। যেমন, $(\Atom{\Obj P}{\Obj a} \land \Atom{\Obj Q}{\Obj b})$
বা~$\lexists[\Obj x][\Atom{\Obj P}{\Obj x}]$-এর মতো অভিব্যক্তি গঠন করা
যায়।

কঠোর অধিযৌক্তিক অধ্যয়নের জন্য প্রয়োজনীয় অর্থতত্ত্বের সংজ্ঞা ও
!!{derivation} পদ্ধতির বিধিগুলি দিতে হলে প্রথমেই নির্ভুলভাবে সংজ্ঞায়িত
করতে হবে, প্রথম-ক্রমের যুক্তিবিদ্যার !!a{sentence} হিসেবে কী গণ্য হয়।
মূল ধারণাটি বোঝা বেশ সহজ: কেবল !!{predicate}s ও !!{constant}s থেকে
$\Atom{\Obj P}{\Obj a}$-এর মতো কয়েকটি সরল !!{sentence}s গঠন করা যায়।
তারপর সংযোজক ও পরিমাণসূচক ব্যবহার করে এগুলি থেকে আরও জটিল বাক্য গঠন করি।
কিন্তু ঠিক কোন বিধিতে আরও জটিল !!{sentence}s গঠন করা অনুমোদিত? সেগুলি
নির্দিষ্ট না করলে “প্রথম-ক্রমের যুক্তিবিদ্যার !!{sentence}”-কে যথেষ্ট
নির্ভুলভাবে সংজ্ঞায়িত করা হয় না। এখানে কয়েকটি সমস্যা আছে। প্রথমটি
হলো, ঠিক প্রতীকক্রমগুলিকেই !!{sentence}s হিসেবে গণ্য করা। দ্বিতীয়টি
হলো, এমনভাবে তা করা যাতে \emph{সব} !!{sentence}s সম্বন্ধে গাণিতিক প্রমাণ
দেওয়া যায়। সব শেষে, !!{sentence}s নিয়ে কয়েকটি প্রাথমিক ক্রিয়ারও
নির্ভুল সংজ্ঞা দিতে হবে; যেমন, “$!A$-র প্রতিটি $\Obj x$-এর বদলে~$\Obj b$
বসাও।” অসুবিধা হলো, প্রথম-ক্রমের যুক্তিবিদ্যার পরিমাণসূচক ও !!{variable}s-এর
কারণে কাজটি কীভাবে করতে হবে তা একেবারে স্পষ্ট নয়। যেমন,
$\lexists[\Obj x][\Atom{\Obj P}{\Obj a}]$ কি !!a{sentence} হিসেবে গণ্য
হওয়া উচিত? $\lexists[\Obj x][\lexists[\Obj x][\Atom{\Obj P}{\Obj x}]]$
সম্বন্ধে কী বলা যায়? “$(\Atom{\Obj P}{\Obj x} \land \lexists[\Obj
x][\Atom{\Obj P}{\Obj x}])$-এ $\Obj x$-এর বদলে~$\Obj b$ বসাও”—এর ফল
কী হওয়া উচিত?

\end{document}
